% Lecture notes for the Latent Dynamics workshop.
%
% This directory is intentionally SELF-CONTAINED: catniplab.sty and
% math_preamble.tex are copied in locally so the notes build from a fresh
% clone of this public repo without any external (memmingCommons) path.
% Trade-off: these copies do not track upstream catniplab.sty fixes - re-copy
% from memmingCommons/latex/ if you want updates.

\documentclass[11pt]{article}
\usepackage[lecture]{catniplab}
% to be included in the preamble
% use % to be included in the preamble
% use % to be included in the preamble
% use \input{math_preamble} in the preamble of the main document

% math & font packages
\usepackage{amssymb}
\usepackage{mathtools} % \DeclarePairedDelimiter
\usepackage{bm} % bold greek

% bold vectors for each alphabet vx
\usepackage{forloop}
\newcommand{\defvec}[1]{\expandafter\newcommand\csname v#1\endcsname{{\mathbf{#1}}}}
\newcounter{ct}
\forLoop{1}{26}{ct}{
    \edef\letter{\alph{ct}}
    \expandafter\defvec\letter
}

% captial \vA
\forLoop{1}{26}{ct}{
    \edef\letter{\Alph{ct}}
    \expandafter\defvec\letter
}

% math
\DeclareMathOperator{\diag}{diag}
\DeclareMathOperator*{\tr}{tr} % trace
\DeclareMathOperator*{\var}{var}
\DeclareMathOperator*{\cov}{cov}

\DeclarePairedDelimiter{\abs}{\lvert}{\rvert}
\DeclarePairedDelimiter{\norm}{\lVert}{\rVert}

\newcommand{\dm}[1]{\ensuremath{\mathrm{d}{#1}}} % dx dy dz dmu
%\newcommand{\RN}[2]{\frac{\dm{#1}}{\dm{#2}}} % (Radon-Nikodym) derivative
\newcommand{\PD}[2]{\frac{\partial #1}{\partial #2}} % partial derivative

\newcommand{\field}[1]{\ensuremath{\mathbb{#1}}}
\newcommand{\reals}{\field{R}}
\newcommand{\complex}{\field{C}}
\newcommand{\integers}{\field{Z}}
\newcommand{\naturalNumbers}{\field{N}}
\newcommand{\trp}{{^\top}} % transpose
\newcommand{\identity}{\ensuremath{\mathbb{I}}}
\newcommand{\ones}{\ensuremath{\mathbf{1}}}
\newcommand{\indicator}[1]{\mathbf{1}_{#1}} % indicator function
\newcommand{\stateSpace}{\reals^n}
\newcommand{\DF}{\nabla_{\vx}\vf} %\newcommand{\DF}{\mathcal{D}\vf}
\newcommand{\vDelta}{\mathbf{\Delta}}
\newcommand{\vPsi}{\mathbf{\Psi}}
\newcommand{\vphi}{\bm{\phi}}

% Exercises and solutions: shared machinery (toggle + printer-friendly boxes).
\input{exercise_setup}
 in the preamble of the main document

% math & font packages
\usepackage{amssymb}
\usepackage{mathtools} % \DeclarePairedDelimiter
\usepackage{bm} % bold greek

% bold vectors for each alphabet vx
\usepackage{forloop}
\newcommand{\defvec}[1]{\expandafter\newcommand\csname v#1\endcsname{{\mathbf{#1}}}}
\newcounter{ct}
\forLoop{1}{26}{ct}{
    \edef\letter{\alph{ct}}
    \expandafter\defvec\letter
}

% captial \vA
\forLoop{1}{26}{ct}{
    \edef\letter{\Alph{ct}}
    \expandafter\defvec\letter
}

% math
\DeclareMathOperator{\diag}{diag}
\DeclareMathOperator*{\tr}{tr} % trace
\DeclareMathOperator*{\var}{var}
\DeclareMathOperator*{\cov}{cov}

\DeclarePairedDelimiter{\abs}{\lvert}{\rvert}
\DeclarePairedDelimiter{\norm}{\lVert}{\rVert}

\newcommand{\dm}[1]{\ensuremath{\mathrm{d}{#1}}} % dx dy dz dmu
%\newcommand{\RN}[2]{\frac{\dm{#1}}{\dm{#2}}} % (Radon-Nikodym) derivative
\newcommand{\PD}[2]{\frac{\partial #1}{\partial #2}} % partial derivative

\newcommand{\field}[1]{\ensuremath{\mathbb{#1}}}
\newcommand{\reals}{\field{R}}
\newcommand{\complex}{\field{C}}
\newcommand{\integers}{\field{Z}}
\newcommand{\naturalNumbers}{\field{N}}
\newcommand{\trp}{{^\top}} % transpose
\newcommand{\identity}{\ensuremath{\mathbb{I}}}
\newcommand{\ones}{\ensuremath{\mathbf{1}}}
\newcommand{\indicator}[1]{\mathbf{1}_{#1}} % indicator function
\newcommand{\stateSpace}{\reals^n}
\newcommand{\DF}{\nabla_{\vx}\vf} %\newcommand{\DF}{\mathcal{D}\vf}
\newcommand{\vDelta}{\mathbf{\Delta}}
\newcommand{\vPsi}{\mathbf{\Psi}}
\newcommand{\vphi}{\bm{\phi}}

% Exercises and solutions: shared machinery (toggle + printer-friendly boxes).
% Shared exercise/solution machinery, \input by both latent_dynamics_notes.tex
% (via math_preamble.tex) and kalman_filter_exercise.tex. Kept in one file so
% the toggle and box styling are not duplicated.
%
% Single source, student/instructor toggle:
%   Student PDF (default): solution boxes are dropped entirely.
%   Instructor PDF: built with \def\SolutionsBuild{} prepended (see the
%   Makefile *_solutions targets). Do NOT hand-edit the toggle here.
% Boxes are printer-friendly: light tint fills (<= 12%) that stay legible in
% grayscale. Colors are defined here so the file is self-contained.

\usepackage{environ}   % capture-or-discard env bodies (robust inside lists)
\usepackage[most]{tcolorbox}
\definecolor{solnframe}{rgb}{0.10,0.25,0.55} % calm blue for solution boxes
\newif\ifsolutions
\ifdefined\SolutionsBuild\solutionstrue\fi

% Inline micro-exercise, numbered per section, shown in both builds. Optional
% argument tags it, e.g. \begin{exercise}[optional stretch] for the harder tier.
\newtcolorbox[auto counter, number within=section]{exercise}[1][]{%
    breakable, enhanced, sharp corners,
    colback=black!3, colframe=black!55, boxrule=0.5pt,
    left=6pt, right=6pt, top=3pt, bottom=3pt,
    coltitle=black, colbacktitle=black!10, fonttitle=\bfseries,
    title={Exercise~\thetcbcounter\if\relax\detokenize{#1}\relax\else\ (#1)\fi},
}

% Solution box, instructor build only. In the student build the environment
% body is captured by environ and simply not emitted (discarded, not scanned),
% which is robust even inside enumerate/itemize.
\newtcolorbox{solnbox}{%
    breakable, enhanced, sharp corners,
    colback=solnframe!5, colframe=solnframe!70, boxrule=0.5pt,
    left=6pt, right=6pt, top=3pt, bottom=3pt,
    coltitle=solnframe!85!black, colbacktitle=solnframe!12,
    fonttitle=\bfseries, title={Solution},
}
\ifsolutions
    \NewEnviron{solution}{\begin{solnbox}\BODY\end{solnbox}}
\else
    \NewEnviron{solution}{}% discard body in the student build
\fi

 in the preamble of the main document

% math & font packages
\usepackage{amssymb}
\usepackage{mathtools} % \DeclarePairedDelimiter
\usepackage{bm} % bold greek

% bold vectors for each alphabet vx
\usepackage{forloop}
\newcommand{\defvec}[1]{\expandafter\newcommand\csname v#1\endcsname{{\mathbf{#1}}}}
\newcounter{ct}
\forLoop{1}{26}{ct}{
    \edef\letter{\alph{ct}}
    \expandafter\defvec\letter
}

% captial \vA
\forLoop{1}{26}{ct}{
    \edef\letter{\Alph{ct}}
    \expandafter\defvec\letter
}

% math
\DeclareMathOperator{\diag}{diag}
\DeclareMathOperator*{\tr}{tr} % trace
\DeclareMathOperator*{\var}{var}
\DeclareMathOperator*{\cov}{cov}

\DeclarePairedDelimiter{\abs}{\lvert}{\rvert}
\DeclarePairedDelimiter{\norm}{\lVert}{\rVert}

\newcommand{\dm}[1]{\ensuremath{\mathrm{d}{#1}}} % dx dy dz dmu
%\newcommand{\RN}[2]{\frac{\dm{#1}}{\dm{#2}}} % (Radon-Nikodym) derivative
\newcommand{\PD}[2]{\frac{\partial #1}{\partial #2}} % partial derivative

\newcommand{\field}[1]{\ensuremath{\mathbb{#1}}}
\newcommand{\reals}{\field{R}}
\newcommand{\complex}{\field{C}}
\newcommand{\integers}{\field{Z}}
\newcommand{\naturalNumbers}{\field{N}}
\newcommand{\trp}{{^\top}} % transpose
\newcommand{\identity}{\ensuremath{\mathbb{I}}}
\newcommand{\ones}{\ensuremath{\mathbf{1}}}
\newcommand{\indicator}[1]{\mathbf{1}_{#1}} % indicator function
\newcommand{\stateSpace}{\reals^n}
\newcommand{\DF}{\nabla_{\vx}\vf} %\newcommand{\DF}{\mathcal{D}\vf}
\newcommand{\vDelta}{\mathbf{\Delta}}
\newcommand{\vPsi}{\mathbf{\Psi}}
\newcommand{\vphi}{\bm{\phi}}

% Exercises and solutions: shared machinery (toggle + printer-friendly boxes).
% Shared exercise/solution machinery, \input by both latent_dynamics_notes.tex
% (via math_preamble.tex) and kalman_filter_exercise.tex. Kept in one file so
% the toggle and box styling are not duplicated.
%
% Single source, student/instructor toggle:
%   Student PDF (default): solution boxes are dropped entirely.
%   Instructor PDF: built with \def\SolutionsBuild{} prepended (see the
%   Makefile *_solutions targets). Do NOT hand-edit the toggle here.
% Boxes are printer-friendly: light tint fills (<= 12%) that stay legible in
% grayscale. Colors are defined here so the file is self-contained.

\usepackage{environ}   % capture-or-discard env bodies (robust inside lists)
\usepackage[most]{tcolorbox}
\definecolor{solnframe}{rgb}{0.10,0.25,0.55} % calm blue for solution boxes
\newif\ifsolutions
\ifdefined\SolutionsBuild\solutionstrue\fi

% Inline micro-exercise, numbered per section, shown in both builds. Optional
% argument tags it, e.g. \begin{exercise}[optional stretch] for the harder tier.
\newtcolorbox[auto counter, number within=section]{exercise}[1][]{%
    breakable, enhanced, sharp corners,
    colback=black!3, colframe=black!55, boxrule=0.5pt,
    left=6pt, right=6pt, top=3pt, bottom=3pt,
    coltitle=black, colbacktitle=black!10, fonttitle=\bfseries,
    title={Exercise~\thetcbcounter\if\relax\detokenize{#1}\relax\else\ (#1)\fi},
}

% Solution box, instructor build only. In the student build the environment
% body is captured by environ and simply not emitted (discarded, not scanned),
% which is robust even inside enumerate/itemize.
\newtcolorbox{solnbox}{%
    breakable, enhanced, sharp corners,
    colback=solnframe!5, colframe=solnframe!70, boxrule=0.5pt,
    left=6pt, right=6pt, top=3pt, bottom=3pt,
    coltitle=solnframe!85!black, colbacktitle=solnframe!12,
    fonttitle=\bfseries, title={Solution},
}
\ifsolutions
    \NewEnviron{solution}{\begin{solnbox}\BODY\end{solnbox}}
\else
    \NewEnviron{solution}{}% discard body in the student build
\fi



\usepackage{cleveref}
\usepackage{amsthm}
\theoremstyle{plain}
\newtheorem{theorem}{Theorem}[section]
\newtheorem{proposition}[theorem]{Proposition}
\newtheorem{corollary}[theorem]{Corollary}
\newtheorem{lemma}[theorem]{Lemma}
\theoremstyle{definition}
\newtheorem{definition}{Definition}[section]
\theoremstyle{remark}
\newtheorem{remark}{Remark}[section]
\theoremstyle{definition}
\newtheorem{funfact}{Fun fact}
\crefname{funfact}{fun fact}{fun facts}
\Crefname{funfact}{Fun fact}{Fun facts}

% local notation (math_preamble already provides \vx \vy \cov \trp \vPsi \identity)
\newcommand{\E}{\mathbb{E}}
\newcommand{\veta}{\bm{\eta}}
\newcommand{\Normal}{\mathcal{N}}
\DeclareMathOperator{\rank}{rank}
\DeclareMathOperator{\col}{col}

% schematic matrix rectangles with dimension labels (cf. the linear algebra notes)
\newcommand*{\clapp}[1]{\hbox to 0pt{\hss#1\hss}}
\newcommand*{\subdims}[3]{\clapp{\raisebox{#1}[0pt][0pt]{$\scriptstyle(#2 \times #3)$}}}
\setlength{\fboxrule}{1pt}

\addbibresource{references.bib}

\title{Latent Dynamics: Lecture Notes}
\author{I.~Memming Park}
\date{\today}

\begin{document}
\maketitle

\begin{abstract}
Notes on two classical linear-Gaussian latent variable models---probabilistic
PCA \parencite{Tipping1999-cv} and factor analysis---with the Gaussian identities
their derivations rely on. Throughout, $\vy \in \reals^m$ is the observed vector, $\vx \in \reals^d$
the latent vector ($m \geq d$, typically $d \ll m$), and $C \in \reals^{m
\times d}$ the loading matrix.
\end{abstract}

\paragraph{How to read these notes.}
The sections are roughly self-contained, so skip any whose results you already
know---most readers can jump straight to probabilistic PCA
(\cref{sec:ppca}). The notation and Gaussian fun facts below are a quick
reference to consult as needed; each fun fact comes with its one-dimensional
special case for intuition. These notes complement the standard textbook
treatment in Bishop's \emph{Pattern Recognition and Machine Learning}
(ch.~12)~\parencite{Bishop2006-to}, which also covers the fully Bayesian version of
PCA.

\section{Notation, from scratch}\label{sec:notation}

A \emph{vector} $\vx = (x_1, \dots, x_d)$ is just a list of $d$ numbers; we write
$\vx \in \reals^d$. Bold lowercase letters ($\vx, \vy, \veta$) are vectors,
capital letters ($C, \Sigma$) are matrices, and $A\trp$ is the transpose of $A$.

The star of these notes is the \emph{Gaussian} (a.k.a.\ \emph{normal})
distribution---the bell curve. In one dimension, a random number $x$ with mean
$\mu$ and variance $\sigma^2$ has the familiar density
\begin{equation}\label{eq:gauss1d}
    p(x) = \frac{1}{\sqrt{2\pi\sigma^2}}
           \exp\!\left(-\frac{(x-\mu)^2}{2\sigma^2}\right).
\end{equation}
In $d$ dimensions the variance $\sigma^2$ is replaced by a $d \times d$
\emph{covariance matrix} $\Sigma$ (it records the variance of each coordinate and
the covariance between every pair), and the mean $\mu$ by a mean \emph{vector}
$\vm$:
\begin{equation}\label{eq:gaussmv}
    p(\vx) = \frac{1}{\sqrt{(2\pi)^d \det \Sigma}}
             \exp\!\left(-\tfrac{1}{2}(\vx-\vm)\trp \Sigma^{-1} (\vx-\vm)\right).
\end{equation}
Writing that mouthful every time is painful, so we abbreviate ``$\vx$ is Gaussian
with mean $\vm$ and covariance $\Sigma$'' as
\begin{equation}\label{eq:normalnotation}
    \vx \sim \Normal(\vm, \Sigma).
\end{equation}
The curly $\Normal$ is just shorthand for \textbf{N}ormal (the distribution is
named after Gauss); the ``$\sim$'' reads ``is distributed as.'' That is the only
new symbol you need---everything below is a statement about what happens to
$\vm$ and $\Sigma$.

\section{Gaussian fun facts}\label{sec:facts}

Both PPCA and factor analysis are \emph{linear-Gaussian}: the observation is a
linear image of a Gaussian latent, plus independent Gaussian noise. Remarkably,
two facts do all the work---how a Gaussian survives an affine map, and how a
jointly Gaussian pair marginalizes and conditions. Gaussians are the friendly
distribution: push them through linear algebra and they stay Gaussian.

\begin{funfact}[Affine maps keep Gaussians Gaussian]\label{fact:affine}
Let $\vx \sim \Normal(\vm, \Sigma)$ and $\vy = A\vx + \vb$ for a fixed matrix
$A$ and vector $\vb$. Then
\begin{equation}\label{eq:affine}
    \vy \sim \Normal(A\vm + \vb,\; A \Sigma A\trp).
\end{equation}
\end{funfact}

\begin{proof}
An affine map of a Gaussian is Gaussian, so it suffices to match the first two
moments. Linearity of expectation gives $\E[\vy] = A\E[\vx] + \vb = A\vm + \vb$.
For the covariance, write $\vy - \E[\vy] = A(\vx - \vm)$, hence
\begin{equation}
    \cov(\vy) = \E\!\left[A(\vx-\vm)(\vx-\vm)\trp A\trp\right]
              = A\,\E\!\left[(\vx-\vm)(\vx-\vm)\trp\right] A\trp
              = A \Sigma A\trp. \qedhere
\end{equation}
\end{proof}

\emph{One dimension.} With scalars, $y = ax + b$ and $x \sim \Normal(\mu,
\sigma^2)$ give $y \sim \Normal(a\mu + b,\; a^2\sigma^2)$: shifting moves the
center, scaling by $a$ stretches the spread by $a^2$.

\begin{funfact}[Independent Gaussians add up]\label{fact:sum}
Let $\vx \sim \Normal(\vm_1, \Sigma_1)$ and $\vz \sim \Normal(\vm_2, \Sigma_2)$
be \emph{independent}. Then their sum is again Gaussian,
\begin{equation}\label{eq:sum}
    \vx + \vz \sim \Normal(\vm_1 + \vm_2,\; \Sigma_1 + \Sigma_2).
\end{equation}
\end{funfact}

\begin{proof}
A sum of independent Gaussians is Gaussian (stack them into
$\begin{psmallmatrix}\vx\\\vz\end{psmallmatrix}$, which is jointly Gaussian, and
apply \cref{fact:affine} with $A = [\,I \;\; I\,]$). Matching moments,
$\E[\vx+\vz] = \vm_1 + \vm_2$, and by independence the cross terms vanish,
\begin{equation}
    \cov(\vx + \vz) = \cov(\vx) + \cov(\vz) = \Sigma_1 + \Sigma_2. \qedhere
\end{equation}
\end{proof}

\emph{One dimension.} Two independent bell curves add to a wider bell:
$\Normal(\mu_1, \sigma_1^2)$ and $\Normal(\mu_2, \sigma_2^2)$ sum to
$\Normal(\mu_1 + \mu_2,\; \sigma_1^2 + \sigma_2^2)$---variances add, so noise
piles up.

Now suppose two vectors are Gaussian \emph{together}. Partition the stacked
vector and its mean and covariance into blocks,
\begin{equation}\label{eq:joint}
    \begin{pmatrix} \va \\ \vb \end{pmatrix}
    \sim \Normal\!\left(
        \begin{pmatrix} \vm_a \\ \vm_b \end{pmatrix},\;
        \begin{pmatrix}
            \Sigma_{aa} & \Sigma_{ab} \\
            \Sigma_{ba} & \Sigma_{bb}
        \end{pmatrix}
    \right),
    \qquad \Sigma_{ba} = \Sigma_{ab}\trp .
\end{equation}

\begin{funfact}[Marginalize and condition]\label{fact:cond}
For the jointly Gaussian pair \eqref{eq:joint}, the \emph{marginal} of $\va$ is
read straight off the blocks,
\begin{equation}\label{eq:marginal}
    \va \sim \Normal(\vm_a,\; \Sigma_{aa}),
\end{equation}
and the \emph{conditional} of $\va$ given $\vb$ is Gaussian with
\begin{align}
    \E[\va \mid \vb] &= \vm_a + \Sigma_{ab}\Sigma_{bb}^{-1}(\vb - \vm_b),
        \label{eq:condmean}\\
    \cov(\va \mid \vb) &= \Sigma_{aa} - \Sigma_{ab}\Sigma_{bb}^{-1}\Sigma_{ba}.
        \label{eq:condcov}
\end{align}
\end{funfact}

\begin{proof}
The marginal is immediate: $\va = [\,I \;\; 0\,]\begin{psmallmatrix}\va\\\vb\end{psmallmatrix}$
is an affine map of a Gaussian, so \cref{fact:affine} gives
$\va \sim \Normal(\vm_a, \Sigma_{aa})$.

For the conditional, subtract off the part of $\va$ that $\vb$ can explain and
define the \emph{residual}
\begin{equation}\label{eq:residual}
    \vr = (\va - \vm_a) - \Sigma_{ab}\Sigma_{bb}^{-1}(\vb - \vm_b).
\end{equation}
Being an affine function of the joint Gaussian \eqref{eq:joint}, $\vr$ is Gaussian
with mean $\mathbf 0$. It is \emph{uncorrelated} with $\vb$,
\begin{equation}
    \cov(\vr, \vb)
        = \Sigma_{ab} - \Sigma_{ab}\Sigma_{bb}^{-1}\Sigma_{bb}
        = 0,
\end{equation}
and for jointly Gaussian vectors uncorrelated means independent. Its covariance is
\begin{equation}
    \cov(\vr)
        = \Sigma_{aa}
          - 2\,\Sigma_{ab}\Sigma_{bb}^{-1}\Sigma_{ba}
          + \Sigma_{ab}\Sigma_{bb}^{-1}\Sigma_{bb}\Sigma_{bb}^{-1}\Sigma_{ba}
        = \Sigma_{aa} - \Sigma_{ab}\Sigma_{bb}^{-1}\Sigma_{ba}.
\end{equation}
Rearranging \eqref{eq:residual},
$\va = \vm_a + \Sigma_{ab}\Sigma_{bb}^{-1}(\vb - \vm_b) + \vr$. Conditioning on
$\vb$ freezes the middle term, while $\vr$---independent of $\vb$---keeps its
distribution. This gives the mean \eqref{eq:condmean} and covariance
\eqref{eq:condcov}.
\end{proof}

Equation \eqref{eq:marginal} says marginalizing is \emph{free}---ignore the
blocks you do not care about. The conditional mean \eqref{eq:condmean} is the
prior mean $\vm_a$ nudged by the observed discrepancy $\vb - \vm_b$, and the
conditional covariance \eqref{eq:condcov} (a \emph{Schur complement}) is smaller
than $\Sigma_{aa}$---observing $\vb$ can only reduce our uncertainty about
$\va$.

\emph{One dimension.} For two correlated scalars with correlation $\rho$, the
conditional mean $\E[a \mid b] = \mu_a + \frac{\sigma_{ab}}{\sigma_b^2}(b -
\mu_b)$ is exactly the least-squares regression line, and the conditional
variance $\sigma_a^2 - \sigma_{ab}^2/\sigma_b^2 = \sigma_a^2(1 - \rho^2)$ shrinks
by the factor $1 - \rho^2$: the more correlated $b$ is with $a$, the more it
tells us.

\section{Probabilistic PCA}\label{sec:ppca}

Probabilistic PCA \parencite{Tipping1999-cv} explains the observed vector as a
low-dimensional latent cause seen through a linear map, corrupted by isotropic
noise. Assume the data has been centered (the sample mean subtracted), so all
means are zero. The \emph{generative model} is
\begin{align}
    \vx &\sim \Normal(\mathbf 0, I_d), \label{eq:ppca-latent}\\
    \vy &= C\vx + \veta, \qquad \veta \sim \Normal(\mathbf 0, \sigma^2 I_m),
        \label{eq:ppca-obs}
\end{align}
with $\vx$ and $\veta$ independent. The latent $\vx \in \reals^d$ is a handful of
hidden factors; the loading matrix $C \in \reals^{m\times d}$ maps them into the
$m$ observed dimensions; and every observed coordinate gets the \emph{same}
noise variance $\sigma^2$---the ``isotropic'' assumption that distinguishes
PPCA from factor analysis later.

Writing out the shapes makes the dimensionality reduction visible. The model
only requires $m \geq d$, but the regime of interest is $d \ll m$: a handful of
latent factors driving many observed coordinates. The loading $C$ is then tall
and thin,
\begin{equation}\label{eq:ppca-shapes}
    \underbrace{\begin{bmatrix} y_1 \\ y_2 \\ \vdots \\ y_m \end{bmatrix}}_{m \times 1}
    =
    \underbrace{\begin{bmatrix}
        C_{11} & \cdots & C_{1d} \\
        C_{21} & \cdots & C_{2d} \\
        \vdots & \ddots & \vdots \\
        C_{m1} & \cdots & C_{md}
    \end{bmatrix}}_{m \times d}
    \underbrace{\begin{bmatrix} x_1 \\ \vdots \\ x_d \end{bmatrix}}_{d \times 1}
    +
    \underbrace{\begin{bmatrix} \eta_1 \\ \eta_2 \\ \vdots \\ \eta_m \end{bmatrix}}_{m \times 1} .
\end{equation}
Reading $C$ by its columns $C = [\,\vc_1 \;\cdots\; \vc_d\,]$, the model is
$\vy = \sum_{k=1}^d x_k \vc_k + \veta$: each latent coordinate $x_k$ turns on a
loading direction $\vc_k \in \reals^m$, and the observation is their weighted sum
plus noise.

\paragraph{What does $\vy$ look like?}
Because $\vy$ is a linear image of the Gaussian $\vx$ plus independent Gaussian
noise, our fun facts hand us its distribution immediately. It is zero-mean, and
combining the affine map (\cref{fact:affine}) on $C\vx$ with the independent sum
(\cref{fact:sum}) for the noise,
\begin{equation}\label{eq:ppca-marginal}
    \vy \sim \Normal\!\left(\mathbf 0,\; \underbrace{C C\trp + \sigma^2 I_m}_{\displaystyle \Sigma_{\vy}}\right).
\end{equation}
So PPCA models the data covariance as a low-rank part $C C\trp$ plus an isotropic
floor $\sigma^2 I_m$. Writing $C = [\,\vc_1\ \cdots\ \vc_d\,]$ through its columns
$\vc_k \in \reals^m$, the low-rank part is a sum of $d$ rank-one outer
products---the tall $m \times d$ block $C$ times the wide $d \times m$ block
$C\trp$,
\begin{equation}\label{eq:ppca-lowrank}
    C C\trp \quad=\quad
    \framebox[0.9cm]{\raisebox{0pt}[1.0cm][1.0cm]{$C$}}
    \;\;
    \framebox[1.8cm]{\raisebox{0mm}[4.5mm][2.5mm]{$C\trp$}}
    \quad=\quad \sum_{k=1}^d \vc_k \vc_k\trp ,
\end{equation}
one bump of shared variance per latent factor. Each $\vc_k \vc_k\trp$ is rank
one, and summing the $d$ of them gives a rank-$d$ matrix ($\rank C C\trp = d$
when the loading columns are linearly independent, and at most $d$ in general).
The observation covariance is these few strong directions of shared variation
riding on a uniform background of noise $\sigma^2 I_m$.

\subsection{Inference: the posterior over the latent}

Given an observation $\vy$, what latent factors likely produced it? We want
$p(\vx \mid \vy)$. The pair $(\vx, \vy)$ is jointly Gaussian---both are linear
functions of the independent Gaussians $\vx$ and $\veta$---so \cref{fact:cond}
applies once we have the covariance blocks. From \eqref{eq:ppca-latent} and
\eqref{eq:ppca-obs},
\begin{equation}
    \cov(\vx) = I_d, \qquad
    \cov(\vy) = C C\trp + \sigma^2 I_m, \qquad
    \cov(\vx, \vy) = \E[\vx(C\vx + \veta)\trp] = C\trp .
\end{equation}
Feeding these into the conditioning formula (\cref{fact:cond}, with $\va = \vx$
and $\vb = \vy$),
\begin{align}
    \E[\vx \mid \vy] &= C\trp(C C\trp + \sigma^2 I_m)^{-1}\vy, \label{eq:ppca-postmean-m}\\
    \cov(\vx \mid \vy) &= I_d - C\trp(C C\trp + \sigma^2 I_m)^{-1} C. \label{eq:ppca-postcov-m}
\end{align}
These invert an $m \times m$ matrix, but with $d \ll m$ we would rather invert
something $d \times d$. The matrix inversion lemma (``push-through identity''),
$(C C\trp + \sigma^2 I_m)^{-1} C = C\,(C\trp C + \sigma^2 I_d)^{-1}$, lets us do
exactly that. Defining the $d \times d$ matrix
\begin{equation}\label{eq:ppca-M}
    M = C\trp C + \sigma^2 I_d,
\end{equation}
the posterior collapses to
\begin{equation}\label{eq:ppca-posterior}
    p(\vx \mid \vy) = \Normal\!\left(M^{-1} C\trp \vy,\; \sigma^2 M^{-1}\right).
\end{equation}
(The covariance follows from
$I_d - C\trp(CC\trp+\sigma^2 I_m)^{-1}C = I_d - M^{-1} C\trp C = I_d - M^{-1}(M - \sigma^2 I_d) = \sigma^2 M^{-1}$,
using $C\trp C = M - \sigma^2 I_d$.) Now the only inverse is the small $d \times d$
matrix $M$. The posterior mean is a linear read-out of $\vy$---not the naive
$C\trp\vy$, but shrunk through $M^{-1}$; as $\sigma^2 \to 0$ it approaches the
least-squares projection $(C\trp C)^{-1}C\trp\vy$ onto the latent subspace.

\subsection{Learning: maximum likelihood}

The parameters $C$ and $\sigma^2$ are fit by maximizing the likelihood of the
data $\{\vy_n\}_{n=1}^N$ under the marginal \eqref{eq:ppca-marginal}. Each
$\vy_n \sim \Normal(\mathbf 0, \Sigma_{\vy})$ with $\Sigma_{\vy} = C C\trp +
\sigma^2 I_m$, so the log-likelihood is
\begin{equation}\label{eq:ppca-loglik}
    \ell(C, \sigma^2)
      = -\frac{N}{2}\Big[\,m \ln(2\pi) + \ln\det\Sigma_{\vy}
        + \tr\!\big(\Sigma_{\vy}^{-1} S\big)\Big],
    \qquad
    S = \frac{1}{N}\sum_{n=1}^N \vy_n \vy_n\trp .
\end{equation}
It touches the data only through the sample covariance $S$---the very statistic
classical PCA diagonalizes.

Recall the \emph{eigendecomposition} of a symmetric positive-semidefinite matrix:
$S$ can be written $S = U \Lambda U\trp$, where $U = [\,\vu_1\ \cdots\ \vu_m\,]$
is orthogonal (its columns $\vu_i$ are orthonormal \emph{eigenvectors}, the
principal directions) and $\Lambda = \diag(\lambda_1, \dots, \lambda_m)$ holds
the \emph{eigenvalues}, ordered $\lambda_1 \geq \cdots \geq \lambda_m \geq 0$.
Each $\lambda_i$ is the variance of the data along direction $\vu_i$. Collect the
top $d$ into $U_d = [\,\vu_1\ \cdots\ \vu_d\,]$ (the $d$ highest-variance
directions) and $\Lambda_d = \diag(\lambda_1, \dots, \lambda_d)$; this
$U_d$ is exactly what PCA returns.

Setting $\partial\ell/\partial C = 0$ has, remarkably, a closed-form global
maximum \parencite{Tipping1999-cv}:
\begin{equation}\label{eq:ppca-ML}
    C_{\mathrm{ML}} = U_d\,(\Lambda_d - \sigma^2 I_d)^{1/2}\,R,
\end{equation}
with $U_d$ and $\Lambda_d$ the top-$d$ eigenvectors and eigenvalues of $S$ just
defined, and $R$ any $d \times d$ orthogonal matrix. That $R$ is the rotation ambiguity we
already saw: $C \mapsto CR$ leaves $C C\trp$, hence the likelihood, unchanged.
The noise variance is the average of the \emph{discarded} eigenvalues,
\begin{equation}\label{eq:ppca-sigmaML}
    \sigma^2_{\mathrm{ML}} = \frac{1}{m - d}\sum_{j = d+1}^{m} \lambda_j ,
\end{equation}
i.e.\ the leftover variance PPCA cannot explain, spread uniformly. So PPCA
recovers the principal subspace: the columns of $C_{\mathrm{ML}}$ span the
top-$d$ eigenvectors of $S$---exactly PCA's directions---with the residual
variance folded into $\sigma^2$. As $\sigma^2 \to 0$ the loadings collapse onto
the scaled principal axes and PPCA reduces to classical PCA.

\paragraph{PPCA is PCA, made probabilistic.}
The two methods return the \emph{same subspace}. Classical PCA keeps the span of
the top $d$ eigenvectors of $S$, namely $\col(U_d)$. The PPCA loading
\eqref{eq:ppca-ML} is $U_d$ right-multiplied by $(\Lambda_d - \sigma^2 I_d)^{1/2}R$---a
diagonal matrix times an orthogonal one, hence an invertible $d \times d$
matrix (the retained $\lambda_1, \dots, \lambda_d$ exceed the average tail
eigenvalue $\sigma^2$, so the square root is real). Right-multiplying by an
invertible matrix never changes the column space, so
\begin{equation}\label{eq:ppca-samespace}
    \col(C_{\mathrm{ML}}) = \col(U_d) :
\end{equation}
the PPCA loadings span the principal subspace \emph{exactly}. The two methods
differ only in how they coordinatize that subspace: PCA scales axis $k$ by
$\sqrt{\lambda_k}$ and pins the orientation, whereas PPCA scales by
$\sqrt{\lambda_k - \sigma^2}$ (shrinking each axis by the noise floor it strips
off) and leaves the orientation free through $R$. As $\sigma^2 \to 0$ even the
scalings coincide, and the posterior mean $\E[\vx \mid \vy]$ becomes the
orthogonal projection of $\vy$ onto $\col(U_d)$---the PCA scores. In short, PPCA
is a probabilistic model whose maximum-likelihood fit \emph{is} PCA: it recasts
what PCA computes as inference in a linear-Gaussian latent variable model.

\paragraph{Two caveats, and the road to factor analysis.}
The equivalence cuts both ways, with a caveat in each direction. First, PCA is
\emph{more general} than its probabilistic reading: as a variance-maximizing
(equivalently, best least-squares subspace) procedure it can be run on data that
are neither Gaussian nor generated by a linear model---PPCA is one probabilistic
interpretation, not a prerequisite. Second, and pulling the other way, PPCA's
isotropic noise $\sigma^2 I_m$ is a strong assumption: every observed coordinate is
forced to carry the \emph{same} private noise variance. Real measurements rarely
oblige---different sensors, neurons, or features have different noise scales.
Relaxing $\sigma^2 I_m$ to a general diagonal noise covariance is exactly factor
analysis, the subject of the next section.

\printbibliography

\end{document}
